\documentclass[tikz,border=10pt]{standalone}
\usepackage{ctex}
\usepackage{amsmath}
\usetikzlibrary{calc, arrows.meta, 3d}

\begin{document}
% 折叠立体示例：平面三角形沿公共边折叠成三棱锥
\begin{tikzpicture}[
  font=\small,
  x={(1.0cm,0cm)},
  y={(0.45cm,0.32cm)},
  z={(0cm,1.0cm)},
  dot/.style={fill=black, circle, inner sep=1.8pt},
  thick,
]

%─── 左侧：展开的平面图 ─────────────────────────────
\begin{scope}[shift={(-5.5, 0, 0)}]
  \node[font=\bfseries, blue!70] at (1.0, 3.0, 0) {折叠前（平面图）};

  % 底面三角形 ABC
  \coordinate (A0) at (0,0,0);
  \coordinate (B0) at (3,0,0);
  \coordinate (C0) at (1.5, 2.2, 0);

  \fill[blue!10, draw=blue!60, thick]
    (A0) -- (B0) -- (C0) -- cycle;
  \node[dot] at (A0) {}; \node[below left, font=\small] at (A0) {$A$};
  \node[dot] at (B0) {}; \node[below right, font=\small] at (B0) {$B$};
  \node[dot] at (C0) {}; \node[above, font=\small] at (C0) {$C$};

  % 将折叠上去的面（△ ACD'）用虚线画，颜色不同
  % D' 在平面内，折叠后会升起
  \coordinate (D0) at (0.0, 4.5, 0);  % 展开时 D 的位置
  \fill[orange!15, draw=orange!70, thick, dashed]
    (A0) -- (C0) -- (D0) -- cycle;
  \node[dot, fill=orange!80] at (D0) {};
  \node[left, font=\small, orange!80!black] at (D0) {$D$（折叠前）};

  % 折叠提示箭头
  \draw[-Stealth, red!70, thick, bend left=25]
    ($(D0)+(0.2,0.3,0)$) to node[right, font=\scriptsize, red]
    {沿 $AC$ 折叠\\$\uparrow$} ($(C0)+(-0.1,0.8,0)$);
\end{scope}

%─── 右侧：折叠后的三棱锥 ──────────────────────────
\begin{scope}[shift={(0.5, 0, 0)}]
  \node[font=\bfseries, orange!80!black] at (1.5, 3.8, 0) {折叠后（三棱锥）};

  % 三棱锥顶点坐标
  \coordinate (A) at (0, 0, 0);
  \coordinate (B) at (3, 0, 0);
  \coordinate (C) at (1.5, 1.8, 0);
  \coordinate (D) at (0.2, 0.8, 2.2);  % D 被折起来了

  % 底面 △ABC（蓝色）
  \fill[blue!12, draw=blue!60, fill opacity=0.8]
    (A) -- (B) -- (C) -- cycle;

  % 侧面（橙色）
  \fill[orange!20, draw=orange!60, fill opacity=0.7]
    (A) -- (C) -- (D) -- cycle;
  \fill[orange!10, draw=orange!40, fill opacity=0.6, dashed]
    (A) -- (B) -- (D) -- cycle;
  \fill[orange!15, draw=orange!50, fill opacity=0.7]
    (B) -- (C) -- (D) -- cycle;

  % 隐藏棱（虚线）
  \draw[gray!60, dashed] (A) -- (D);

  % 顶点
  \node[dot] at (A) {}; \node[below left, font=\small] at (A) {$A$};
  \node[dot] at (B) {}; \node[below right, font=\small] at (B) {$B$};
  \node[dot] at (C) {}; \node[right, font=\small] at (C) {$C$};
  \node[dot, fill=orange!80] at (D) {}; \node[above, font=\small, orange!80!black] at (D) {$D$};

  % 折叠棱 AC 标注
  \draw[red, very thick] (A) -- (C)
    node[midway, left, font=\scriptsize, red] {折叠轴 $AC$};

  % 折叠角标注
  \draw[<->, red!60, thin]
    ($(A)+(0.15,0.05,0.9)$) -- ($(A)+(0.05,0.55,0)$)
    node[midway, right, font=\scriptsize, red] {折叠角 $\alpha$};
\end{scope}

%─── 底部说明 ─────────────────────────────────────
\node[font=\scriptsize, gray, align=center] at (2.0, -1.2, 0)
  {折叠时，原共面点变为不共面；原共面角（如 $\angle ACD$）不变，但位置关系改变。\\
   关键：折叠保持棱长不变；$D$ 到底面 $ABC$ 的距离由折叠角决定。};

\end{tikzpicture}
\end{document}
